% !PRINTED 61
\weberpage{61}
%<<CONT
\[
\begin{aligned}
A &\equiv -1 \pmod{2^{\lambda}} &&\equiv 1 \pmod{q_1^{\varkappa_1}} &&\equiv 1 \pmod{q_2^{\varkappa_2}} \ldots\\
B &\equiv \phantom{-}5\;'' &&\equiv 1\;'' &&\equiv 1\;'' \ldots\\
C_1 &\equiv \phantom{-}1\;'' &&\equiv g_1\;'' &&\equiv 1\;'' \ldots\\
C_2 &\equiv \phantom{-}1\;'' &&\equiv 1\;'' &&\equiv g_2\;'' \ldots
\end{aligned}
\]

\noindent
$\ldots\ldots\ldots\ldots\ldots\ldots\ldots\ldots\ldots\ldots\ldots\ldots\ldots\ldots\ldots\ldots\ldots\ldots\ldots\ldots\ldots\ldots$

und nach 4. erh\"alt man dann f\"ur jede Zahl $n$ unserer Gruppe

\weq{3}{n \equiv A^{\alpha} B^{\beta} C_1^{\gamma_1} C_2^{\gamma_2} \ldots \pmod{m}.}

Die $A, B, C_1, C_2 \ldots$ sind also die Elemente einer Basis der
Gruppe $\mathfrak{N}$ von den Graden $a, b, c_1, c_2 \ldots$

Wenn $m$ ungerade oder nur durch die erste Potenz von 2
theilbar ist, so fallen aus der Basis die beiden Elemente $A, B$
weg. Ist $m$ durch 4, aber nicht durch 8 theilbar, so f\"allt $B$ weg
und das Element $A$ vom zweiten Grade bleibt.

\gesperrt{5. Zerlegt man die Zahlen $a, b, c_1, c_2 \ldots$ in Potenzen
von einander verschiedener Primzahlen, so erh\"alt man die Invarianten der Gruppe $\mathfrak{N}$.}

Um die verschiedenen F\"alle zusammenzufassen, bezeichnen
wir die Elemente der Basis (2) mit

\weq{4}{C_{-1}, C_0, C_1 \ldots C_\mu.}

Hierin sollen $C_{-1}, C_0$ f\"ur $A$ und $B$ stehen und sind also
gleich 1 zu setzen, wenn $\lambda = 0$ oder $\lambda = 1$ ist. Wenn $\lambda = 2$
ist, so ist $B = C_{-1} = 1$, $A = C_0$, und wenn $\lambda > 1$ ist, $A = C_{-1}$,
$B = C_0$ zu setzen.
% [sic: Weber prints \lambda > 1; the preceding cases require \lambda > 2 here.]

Die Grade dieser Elemente seien mit

\weq{5}{c_{-1}, c_0, c_1 \ldots c_\mu}

bezeichnet, und die Indices einer Zahl aus $\mathfrak{N}$, die nach den
Gr\"ossen (5) als Moduln zu nehmen sind, mit

\weq{6}{\nu_{-1}, \nu_0, \nu_1 \ldots \nu_\mu.}

Ist $\lambda = 0$ oder 1, so haben $\nu_{-1}$ und $\nu_0$ nur den Werth 0,
ist $\lambda = 2$, so hat $\nu_{-1}$ den Werth 0 und $\nu_0$ den Werth 0 oder 1.
Ist $\lambda \geqq 3$, so hat $\nu_{-1}$ einen der beiden Werthe 0, 1, und $\nu_0$
einen der Werthe $0, 1, 2, \ldots c_0 - 1$; $\mu$ ist immer gleich der
Anzahl der von einander verschiedenen ungeraden Primzahlen,
die in $m$ aufgehen. Wir wollen, indem wir $\nu_{-1}, c_{-1}$ bei Seite
lassen, $\nu_0, \nu_1 \ldots \nu_\mu$ die den Primzahlen $2, q_1 \ldots q_\mu$ entsprechenden
\gesperrt{Indices} von $n$ und $c_0, c_1 \ldots c_\mu$ die denselben Primzahlen entsprechenden
\gesperrt{Indexmoduln} nennen. Diese Indexmoduln sind
\[
c_0 = \varphi(2^{\lambda-1}), \quad c_1 = \varphi(q_1^{\varkappa_1}), \ldots c_\mu = \varphi(q_\mu^{\varkappa_\mu}),
\]
%CONT>>
